%% ============================================================
%%  chapitre6.tex — الفصل السادس: دور مفتش الشغل ومراقبة احترام القانون
%% ============================================================

\chapter{دور مفتش الشغل ومراقبة احترام القانون}

\section{اختصاصات مفتش الشغل}

يُعدّ مفتش الشغل الذراعَ التنفيذية لتطبيق تشريعات الشغل، ويتمتع باختصاصات واسعة:

\subsection{الاختصاصات الرقابية}
\begin{itemize}[leftmargin=*, itemsep=6pt]
    \item التفتيش على المقاولات والمصانع والمنشآت في أي وقت.
    \item الاطلاع على العقود والسجلات والوثائق الإدارية والاجتماعية.
    \item استجواب المشغّلين والأجراء وجمع الإفادات.
    \item أخذ عيّنات للفحص في مجال الصحة والسلامة.
\end{itemize}

\subsection{الاختصاصات الإصلاحية والتوفيقية}
\begin{itemize}[leftmargin=*, itemsep=6pt]
    \item التدخل لتسوية النزاعات الفردية والجماعية بالتوفيق.
    \item إصدار الإنذارات وإعذار المشغّلين بتصحيح الوضع.
    \item تقديم المشورة والتوضيحات القانونية.
\end{itemize}

%% ─────────────────────────────────────────────────────────
\section{مساطر المراقبة والمعاينة}

تسير عملية التفتيش وفق مساطر منظمة:

\begin{enumerate}[leftmargin=*, itemsep=8pt]
    \item \textbf{الزيارة الميدانية:} يزور مفتش الشغل المنشأة بصفة اعتيادية
          أو إثر شكاية.
    \item \textbf{الاطلاع على الوثائق:} فحص سجل الأجور، قوائم الأجراء،
          عقود العمل، وثائق الضمان الاجتماعي.
    \item \textbf{الاستماع:} الاستماع إلى المشغّل وممثلي الأجراء والأجراء أنفسهم.
    \item \textbf{التحقق من الواقعة:} معاينة المخالفات على عين المكان وتوثيقها.
    \item \textbf{إعطاء مهلة التصحيح:} منح المشغّل مهلة معقولة لتصحيح الوضع
          قبل إحالة المحضر إلى النيابة العامة.
\end{enumerate}

%% ─────────────────────────────────────────────────────────
\section{المحاضر والإنذارات والمتابعات}

\subsection{المحاضر الرسمية}
يُحرّر مفتش الشغل محاضر رسمية موثِّقة للمخالفات المعاينة، وتُرسل نسخ منها إلى:
\begin{itemize}[leftmargin=*, itemsep=4pt]
    \item النيابة العامة.
    \item وزارة التشغيل.
    \item المشغّل المعني.
\end{itemize}

\subsection{الإنذارات}
قبل تحرير المحضر الجنائي، يصدر مفتش الشغل إنذاراً رسمياً للمشغّل يتضمن:
\begin{itemize}[leftmargin=*, itemsep=4pt]
    \item وصف المخالفة.
    \item الأحكام القانونية المخالَفة.
    \item الأجل الممنوح للتصحيح.
    \item التبعات القانونية في حالة عدم التصحيح.
\end{itemize}

%% ─────────────────────────────────────────────────────────
\section{حدود التفتيش والصعوبات العملية}

\begin{tcolorbox}[colback=lightgray, colframe=maincolor, arc=4pt,
    title={\textbf{تحديات ميدان التفتيش}}]
\begin{itemize}[leftmargin=*, itemsep=4pt]
    \item \textbf{قلة الموارد البشرية:} نسبة مفتشي الشغل إلى عدد المقاولات
          لا تزال دون المستوى المطلوب.
    \item \textbf{تعقيد منظومة التعاقد الباطني:} صعوبة كشف حالات المناولة المقنّعة.
    \item \textbf{ضعف التنسيق:} بين مصالح التفتيش والضمان الاجتماعي والقضاء.
    \item \textbf{مقاومة بعض المقاولات:} رفض الوصول أو إخفاء الوثائق.
    \item \textbf{ضعف الوعي القانوني:} لدى الأجراء المؤقتين بحقوقهم.
\end{itemize}
\end{tcolorbox}

